\begin{derivation}[从\eqnrefs{eq:wick-generating-identity-dp74}到\eqnrefs{eq:wick-general-r3}]
把生成恒等式写成“正规序线性源”与“收缩二次源”的乘积：
\begin{equation*}
 \mathcal T\exp\!\left[
 \frac{\ii}{\hbar c}\int\dd^4x\,\mathcal J(x)\varphi(x)
 \right]
 =:\exp\!\left[
 \frac{\ii}{\hbar c}\int\dd^4x\,\mathcal J(x)\varphi(x)
 \right]:
 \exp\!\left[-\frac{1}{2\hbar^2c^2}
 \int\dd^4x\dd^4y\,\mathcal J(x)\Delta_F(x-y)\mathcal J(y)
 \right].
\end{equation*}
要得到 $n$ 个指定场 $\varphi_r\equiv\varphi(x_r)$ 的时间序乘积，对两边作 $n$ 次源泛函微分后令 $\mathcal J=0$。非零项中，每一个微分只有两种去处：命中正规序指数的一次源，留下一个未收缩场；或者与另一个微分共同命中二次源指数中的同一个 $\mathcal J(x)\mathcal J(y)$，产生一条二点收缩 $\Delta_F(x_r-x_s)$。

固定恰有 $k$ 条收缩。需要从 $n$ 个标签中选择 $2k$ 个参与收缩，再把它们配成 $k$ 对。带标签配对数为
\begin{equation*}
 \binom{n}{2k}\frac{(2k)!}{2^k k!}
 =\frac{n!}{(n-2k)!2^k k!}.
\end{equation*}
这个组合因子可以直接从二次指数的展开系数看出。其第 $k$ 阶是
\begin{equation*}
 \frac{1}{k!}
 \left[-\frac{1}{2\hbar^2c^2}
 \int\dd^4x\dd^4y\,
 \mathcal J(x)\Delta_F(x-y)\mathcal J(y)
 \right]^k.
\end{equation*}
$n$ 次微分中命中这 $2k$ 个源的排列产生 $(2k)!$；每条二次因子内部交换两个端点产生 $2^k$ 重数，而 $k$ 条二次因子之间交换又产生 $k!$ 重数。原系数中的 $1/(2^kk!)$ 恰好把这些“同一无序配对的重复计数”除掉，所以每一个具体部分配对只留下系数 1。

因此若 $P$ 表示标签集合 $\{1,\ldots,n\}$ 的任意部分配对，$V(P)$ 表示所有已参与配对的标签，则
\begin{equation*}
 \mathcal T\!\left\{\prod_{r=1}^n\varphi_r\right\}
 =\sum_{P}
 \left[\prod_{(r,s)\in P}\Delta_F(x_r-x_s)\right]
 :\!\prod_{j\notin V(P)}\varphi_j\!:,
\end{equation*}
即正文\eqnrefs{eq:wick-general-r3}。

例如 $n=4$ 时，$k=2$ 的完全收缩数为
$4!/(2^2 2!)=3$，正是
$\Delta_{12}\Delta_{34}$、$\Delta_{13}\Delta_{24}$、$\Delta_{14}\Delta_{23}$ 三项。这个低阶检查同时说明 Wick 定理的组合系数来自源导数，而不是费曼图规则额外规定的权重。
\end{derivation}

\begin{derivation}[从\eqnrefs{eq:phi4-first-dyson-r68}到\eqnrefs{eq:phi4-tree-green-r3}]

正文\eqnrefs{eq:phi4-first-dyson-r68} 是 $\varphi^4$ 理论四点函数的一阶 Dyson 插入。连通树级项要求每个外场都与顶点 $z$ 上四个场中的一个且仅一个收缩，因此共有 $4!$ 种相同配对；它们与作用量中的 $1/4!$ 抵消后得到正文\eqnrefs{eq:phi4-tree-green-r3}。



四点函数的一阶 Dyson 项为
\begin{equation*}
 G^{(4,1)}(x_1,x_2,x_3,x_4)
 =-\frac{\ii g_4}{4!\,\hbar^2c^2}
 \int\dd^4 z\,
 \langle0|\Tord\{\varphi_1\varphi_2\varphi_3\varphi_4\varphi^4(z)\}|0\rangle .
\end{equation*}
要得到“一个顶点连接四个外点”的连通树图，每个外部场必须与 $z$ 点四个场中的一个且仅一个收缩。先给 $\varphi_1$ 选 $z$ 点场有 $4$ 种；再给 $\varphi_2$ 选剩余场有 $3$ 种；之后分别有 $2$ 种和 $1$ 种，所以
\begin{equation*}
 N_{\rm contr}=4\times3\times2\times1=4!.
\end{equation*}
每一种收缩都给出同一个核乘积
\begin{equation*}
 \Delta_F(x_1-z)\Delta_F(x_2-z)\Delta_F(x_3-z)\Delta_F(x_4-z).
\end{equation*}
因此 $4!$ 个相同项与相互作用中的 $1/4!$ 抵消，得到
\begin{equation*}
 G_c^{(4)}
 =-\frac{\ii g_4}{\hbar^2c^2}
 \int\dd^4 z\,\prod_{a=1}^4\Delta_F(x_a-z),
\end{equation*}
即正文\eqnrefs{eq:phi4-tree-green-r3}。

再把四条外线分别 Fourier 展开。四个传播子的指数中所有 $z$ 依赖合成
\begin{equation*}
 \exp\!\left[-\frac{\ii}{\hbar}
 (p_1+p_2-p_3-p_4)\cdot z\right].
\end{equation*}
于是顶角位置积分给
\begin{equation*}
 \int\dd^4 z\,
 \ee^{-\ii(p_1+p_2-p_3-p_4)\cdot z/\hbar}
 =(2\pi\hbar)^4\delta^{(4)}(p_1+p_2-p_3-p_4).
\end{equation*}
所以顶角四动量守恒正是局域相互作用对顶角位置积分后的 Fourier $\delta$ 函数。
\end{derivation}

\begin{derivation}[从\eqnrefs{eq:golden-sinc-r3}到\eqnrefs{eq:fermi-golden-r3}]
记 $\Delta E=E_f-E_i$。一阶振幅为

\begin{align*}
 c_f^{(1)}(T)
 &=-\frac{\ii}{\hbar}V_{fi}
 \int_{-T/2}^{T/2}\dd t\,\ee^{\ii\Delta E t/\hbar}\\
 &=-\frac{\ii}{\hbar}V_{fi}
 \left[\frac{\hbar}{\ii\Delta E}
 \ee^{\ii\Delta E t/\hbar}\right]_{-T/2}^{T/2}\\
 &=-\frac{2\ii V_{fi}}{\Delta E}
 \sin\!\left(\frac{\Delta E T}{2\hbar}\right).
\end{align*}
把 $\sinc x\equiv\sin x/x$ 提出，得到
\begin{equation*}
 |c_f^{(1)}(T)|^2
 =\frac{|V_{fi}|^2T^2}{\hbar^2}
 \sinc^2\!\left(\frac{\Delta E T}{2\hbar}\right),
\end{equation*}
即\eqnrefs{eq:golden-sinc-r3}。它在有限 $T$ 时不是 $\delta$ 函数；主峰宽度量级是 $\hbar/T$。

对连续末态求和，概率为
\begin{equation*}
 P_{i\to f}(T)
 =\int\dd E\,\nu_f(E)
 \frac{|V_{fi}(E)|^2T^2}{\hbar^2}
 \sinc^2\!\left[\frac{(E-E_i)T}{2\hbar}\right].
\end{equation*}
这里 $\nu_f(E)$ 是末态每单位能量的态密度，SI 量纲为 $\mathrm{J^{-1}}$。有限时间谱核与 Dirac $\delta$ 分布的对应由下列核固定：
\begin{equation*}
 K_T(\Delta E)
 \equiv
 \frac{T}{2\pi\hbar}
 \sinc^2\!\left(\frac{\Delta E\,T}{2\hbar}\right).
\end{equation*}
它具有能量倒数量纲。作变量代换 $x=\Delta E T/(2\hbar)$ 后，
\begin{align*}
 \int_{-\infty}^{+\infty}\dd(\Delta E)\,K_T(\Delta E)
 &=\frac{T}{2\pi\hbar}
 \frac{2\hbar}{T}
 \int_{-\infty}^{+\infty}\dd x\,\sinc^2x\\
 &=\frac1\pi\int_{-\infty}^{+\infty}\dd x\,\sinc^2x=1.
\end{align*}
因此 $K_T$ 是面积固定为一、宽度随 $T^{-1}$ 收缩的归一化谱核；对任意在 $E_i$ 附近平滑的测试函数 $F(E)$，$T\to\infty$ 时满足
\begin{equation*}
 \int\dd E\,F(E)K_T(E-E_i)\longrightarrow F(E_i),
\end{equation*}
也就是 $K_T(\Delta E)\to\delta(\Delta E)$ 的分布极限。
令
\begin{equation*}
F(E)\equiv |V_{fi}(E)|^2\nu_f(E).
\end{equation*}
由 $K_T$ 的定义，有限时间概率可以严格改写为
\begin{equation*}
P_{i\to f}(T)
=\frac{2\pi T}{\hbar}
\int\dd E\,F(E)K_T(E-E_i).
\end{equation*}
再定义有限时间余项
\begin{equation*}
\delta F_T
\equiv
\int\dd E\,[F(E)-F(E_i)]K_T(E-E_i),
\end{equation*}
利用 $\int K_T\dd E=1$ 得到精确等式
\begin{equation*}
P_{i\to f}(T)
=\frac{2\pi T}{\hbar}[F(E_i)+\delta F_T].
\end{equation*}
前面已经证明 $K_T\to\delta$ 的分布极限，所以对满足该近似恒等核收敛条件的连续 $F$，有 $\lim_{T\to\infty}\delta F_T=0$。因此长时间跃迁率严格由极限定义
\begin{equation*}
\lim_{T\to\infty}\frac{P_{i\to f}(T)}{T}
=\frac{2\pi}{\hbar}F(E_i).
\end{equation*}
也就是
\begin{equation*}
 w_{i\to f}=\frac{2\pi}{\hbar}|V_{fi}|^2\nu_f(E_i),
\end{equation*}
即 Fermi 黄金律~\eqnrefs{eq:fermi-golden-r3}。能量 $\delta$ 分布由有限时间谱核与连续谱卷积的长时间极限产生；有限 $T$ 时应保留具有有限宽度的 $\sinc^2$ 核。
\end{derivation}

% Detailed derivations gathered at the end of the owning structure.

\begin{derivation}[从\eqnrefs{eq:generic-fermion-contraction-dp8}到\eqnrefs{eq:fermion-wick-four-r68}]

正文\eqnrefs{eq:generic-fermion-contraction-dp8} 定义了自由费米场的一次收缩。对两对费米场只有直接配对和交叉配对两种完全收缩，而后者相对前者多一次 Grassmann 奇置换；把两项相加便得到正文\eqnrefs{eq:fermion-wick-four-r68}。



考虑固定初始顺序
\begin{equation*}
 \mathscr W_4
 =\langle0|T\{\psi_1\bar\psi_2\psi_3\bar\psi_4\}|0\rangle.
\end{equation*}
自由真空期望值必须把每个 $\psi$ 与一个 $\bar\psi$ 配对。第一种配对是 $(1,2)(3,4)$；两对在原顺序中已经相邻，因此不需要额外奇置换，贡献
\begin{equation*}
 +S_{12}S_{34}.
\end{equation*}
第二种配对是 $(1,4)(3,2)$。$\bar\psi_4$ 与 $\psi_1$ 收缩时必须越过中间的费米算符；等价地说，从配对 $(1,2)(3,4)$ 变到交叉配对改变了一次配对排列的宇称，因此贡献带负号。于是
\begin{equation*}
 {
 \mathscr W_4=S_{12}S_{34}-S_{14}S_{32}.}
 \end{equation*}
这个负号不是某个 QED 顶角的性质；即使没有电磁相互作用，自由费米高斯也已经具有它。

一般地，考虑 $n$ 个 $\psi_i$ 与 $n$ 个 $\bar\psi_j$。每个完全收缩等价于一个置换 $\pi\in S_n$，把第 $i$ 个 $\psi_i$ 配到 $\bar\psi_{\pi(i)}$。每次交换 Grassmann-odd 算符都产生 $-1$，故
\begin{equation*}
 \langle0|T\{\psi_1\bar\psi_1\cdots\psi_n\bar\psi_n\}|0\rangle
 =\sum_{\pi\in S_n}
 \operatorname{sgn}(\pi)
 \prod_{i=1}^{n}S_{i,\pi(i)}.
 \end{equation*}
右边正是矩阵 $S_{ij}$ 的行列式：
\begin{equation*}
 {
 \langle T\,\psi_1\bar\psi_1\cdots\psi_n\bar\psi_n\rangle_0
 =\det[S_{ij}].}
 \end{equation*}
与此相对，玻色高斯的完全收缩是不含 $\operatorname{sgn}(\pi)$ 的积和式配对和。Grassmann 高斯积分产生行列式，普通高斯积分产生行列式的负幂；它们与这里的 Wick 组合计数是同一反对易结构的不同表现。
\end{derivation}

\begin{derivation}[从\eqnrefs{eq:generic-fermion-contraction-dp8}到\eqnrefs{eq:generic-closed-fermion-loop-r68}]
在相互作用展开中，若若干费米收缩连成闭环，例如取如下示意性的双线性乘积

\begin{equation*}
 (\bar\psi_1\Gamma_1\psi_1)
 (\bar\psi_2\Gamma_2\psi_2)\cdots
 (\bar\psi_n\Gamma_n\psi_n),
\end{equation*}
并取循环收缩
\begin{equation*}
 \psi_1\leftrightarrow\bar\psi_2,
 \quad
 \psi_2\leftrightarrow\bar\psi_3,
 \quad\ldots\quad
 \psi_n\leftrightarrow\bar\psi_1,
\end{equation*}
则配对置换是一个长度为 $n$ 的循环。长度为 $n$ 的循环可分解成 $n-1$ 次对换，所以其宇称为 $(-1)^{n-1}$。另一方面，把每个双线性型中的 $\bar\psi_i$ 移到用于写矩阵迹的统一位置还产生 $(-1)^n$。两者相乘为
\begin{equation*}
 (-1)^{n-1}(-1)^n=-1.
\end{equation*}
因此一个闭合费米子圈相对于同样连接结构的对易的场多出
\begin{equation*}
 {-1.}
 \end{equation*}
下一章 QED 真空极化的迹
$-\operatorname{Tr}[\gamma^\mu S_F\gamma^\nu S_F]$
就是这个一般结论的具体实例。
\end{derivation}

\begin{derivation}[从\eqnrefs{eq:lsz-pole-plus-cont-r7}到\eqnrefs{eq:phi4-lsz-result-r7}]
先只处理一条外腿，因为多腿公式只是把同一个极点提取操作重复 $n$ 次。设 $\zeta=p^2-m_{\rm phys}^2c^2$，并把 $\zeta$ 暂时视为复变量。对时序关联函数中与坐标 $x$ 对应的那一个场插入完整态，稳定单粒子态给出
\begin{equation*}
 \langle\Omega|\varphi_H(0)|p,\lambda\rangle
 =Z^{1/2}\,\mathcal N_p,
\end{equation*}
其中 $\mathcal N_p$ 是由全书协变态归一化固定的运动学因子。把该单粒子贡献与其余所有中间态分开，Fourier 变换后必然具有形式
\begin{equation*}
 \widetilde G_{n,c}(p,\ldots)
 =\frac{\ii\hbar^3c\,Z^{1/2}}
 {\zeta}\,
 \widetilde F_{n-1}(p;\ldots)
 +\widetilde R^{(1)}_n(p,\ldots),
\end{equation*}
这里 $\widetilde F_{n-1}$ 是已经把这一条外腿换成归一化单粒子态后的矩阵元；$\widetilde R^{(1)}_n$ 包含连续谱和在这一外部通道中正则的项。稳定粒子极点与多粒子阈值分离，所以在 $\zeta=0$ 的一个小邻域内
\begin{equation*}
 \lim_{\zeta\to0}\zeta\,\widetilde R^{(1)}_n=0.
\end{equation*}
因此极点留数不是一句“外腿传播子可以删掉”，而是如下复分析极限：
\begin{align*}
 \lim_{\zeta\to0}
 Z^{-1/2}\frac{\zeta}{\ii\hbar^3c}
 \widetilde G_{n,c}
 &=\lim_{\zeta\to0}
 \left[
 \widetilde F_{n-1}
 +Z^{-1/2}\frac{\zeta}{\ii\hbar^3c}
 \widetilde R^{(1)}_n
 \right]\\
 &=\widetilde F_{n-1}(p_{\rm on};\ldots).
\end{align*}
Feynman 的 $+\ii0^+$ 是实动量边界值的处方；截肢本身应理解为先在复 $\zeta$ 平面取极点留数，再回到质量壳，因此不会出现把 $\zeta/(\zeta+\ii0^+)$ 当作普通数值比值的歧义。

对第二、第三直到第 $n$ 条外腿重复同一操作。若定义 $\zeta_i=p_i^2-m_i^2c^2$，则联合极点邻域中的多变量 Laurent 展开可写成
\begin{align*}
 \widetilde G_{n,c}
 ={}&(2\pi\hbar)^4\delta^{(4)}\!\left(\sum_i\sigma_i^{\rm io}p_i\right)
 \left[\prod_i
 \frac{\ii\hbar^3c\,Z_i^{1/2}}{\zeta_i}\right]
 \ii\mathfrak A_{fi}(\{p_i\})
 +\widetilde R_n,\\
 &\lim_{\zeta_1,\ldots,\zeta_n\to0}
 \left(\prod_i\zeta_i\right)\widetilde R_n=0.
\end{align*}
第二行正是“余项至少缺少一条完整外腿极点”的数学含义。于是
\begin{align*}
 &\lim_{\zeta_1,\ldots,\zeta_n\to0}
 \left[\prod_iZ_i^{-1/2}
 \frac{\zeta_i}{\ii\hbar^3c}\right]
 \widetilde G_{n,c}\\
 &\qquad=(2\pi\hbar)^4\delta^{(4)}\!\left(\sum_i\sigma_i^{\rm io}p_i\right)
 \ii\mathfrak A_{fi}(\{p_{i,\rm on}\}),
\end{align*}
这就得到正文的动量空间 LSZ 公式。注意截肢和在壳极限的顺序来自极点留数的定义：先乘 $\zeta_i$ 提取留数，再令 $\zeta_i\to0$；若反过来，原 Green 函数正好落在极点上。

现在把这一抽象结构用于树级 $\varphi^4$。一次 Dyson 插入给
\begin{equation*}
 -\frac{\ii}{\hbar c}
 \int\dd^4x\,
 \frac{g_4}{4!\,\hbar c}\,
 \varphi_I(x)^4.
\end{equation*}
把四个外部场分别与顶角处四个场收缩有 $4!$ 种双射，恰好消掉相互作用中的 $1/4!$。每条收缩给一个自由 Feynman 传播子，所以 Fourier 变换后的连通四点函数为
\begin{equation*}
 \widetilde G^{(1)}_{4,c}
 =-\frac{\ii g_4}{\hbar^2c^2}
 (2\pi\hbar)^4\delta^{(4)}(p_1+p_2-p_3-p_4)
 \prod_{a=1}^{4}
 \frac{\ii\hbar^3c}{p_a^2-m^2c^2+\ii0^+}.
\end{equation*}
在树级 $Z=1$。令 $\zeta_a=p_a^2-m^2c^2$，逐腿乘逆极点因子后，乘积中的四条外腿逐一消去：
\begin{align*}
 &\left[\prod_{a=1}^{4}
 \frac{\zeta_a}{\ii\hbar^3c}\right]
 \widetilde G^{(1)}_{4,c}\\
 &\quad=-\frac{\ii g_4}{\hbar^2c^2}
 (2\pi\hbar)^4\delta^{(4)}(p_1+p_2-p_3-p_4)
 \prod_{a=1}^{4}
 \frac{\zeta_a}{\zeta_a+\ii0^+}.
\end{align*}
按上面说明的极点留数顺序取极限，四个因子都贡献单位留数，于是得到正文\eqnrefs{eq:phi4-lsz-result-r7}。因此树级顶角规则不是另记的一条口诀：$4!$ 来自 Wick 收缩计数，四条传播子来自外部收缩，而 LSZ 的四个逆极点恰好把这些外腿传播子截去，只留下局域顶角系数。
\end{derivation}


\begin{derivation}[从\eqnrefs{eq:lsz-finite-time-mode-projection-dp52}到\eqnrefs{eq:lsz-one-leg-position-dp52}]
先验证有限时刻投影式的归一化。把自由场展开~\eqnrefs{eq:scalar-mode-r2} 与

\begin{equation*}
 \pi(x)=\frac{1}{c^2}\partial_t\varphi(x)
 =-\ii\int\frac{\dd^3 q}{(2\pi\hbar)^3}
 \sqrt{\frac{E_{\mathbf q}}{2c^2}}
 \left[
 a(\mathbf q)\ee^{-\ii q\cdot x/\hbar}
 -a^\dagger(\mathbf q)\ee^{+\ii q\cdot x/\hbar}
 \right]
\end{equation*}
代入\eqnrefs{eq:lsz-finite-time-mode-projection-dp52}。对 $a(\mathbf q)$ 分支，方括号中的两个系数分别给
\begin{align*}
 \frac{1}{\hbar c}\sqrt{\frac{E_{\mathbf p}}{2}}
 \frac{\hbar c}{\sqrt{2E_{\mathbf q}}}
 &\xrightarrow{\mathbf q=\mathbf p}\frac12,\\
 \frac{\ii c}{\sqrt{2E_{\mathbf p}}}
 \left(-\ii\sqrt{\frac{E_{\mathbf q}}{2c^2}}\right)
 &\xrightarrow{\mathbf q=\mathbf p}\frac12.
\end{align*}
空间积分给
\begin{equation*}
 \int\dd^3 x\,
 \ee^{-\ii(\mathbf p-\mathbf q)\cdot\mathbf x/\hbar}
 =(2\pi\hbar)^3\delta^{(3)}(\mathbf p-\mathbf q),
\end{equation*}
因此两个 $1/2$ 相加后恰好得到 $a(\mathbf p)$。对 $a^\dagger(\mathbf q)$ 分支，空间积分要求 $\mathbf q=-\mathbf p$，而两个系数变成 $+1/2$ 与 $-1/2$，所以完全抵消。由此，\eqnrefs{eq:lsz-finite-time-mode-projection-dp52} 对自由场确实是模展开的精确逆变换。

现在把场换成完整 Heisenberg 场，并记
\begin{equation*}
 f_p(x)\equiv \ee^{+\ii p\cdot x/\hbar},
 \qquad p^0=E_{\mathbf p}/c.
\end{equation*}
用 $\pi_H=c^{-2}\partial_t\varphi_H$，\eqnrefs{eq:lsz-finite-time-mode-projection-dp52} 可写成
\begin{equation*}
 a_t(\mathbf p)
 =\int\dd^3 x\,f_p(x)
 \left[
 A_p\varphi_H(x)+B_p\partial_t\varphi_H(x)
 \right],
\end{equation*}
其中
\begin{equation*}
 A_p=\frac{1}{\hbar c}\sqrt{\frac{E_{\mathbf p}}{2}},
 \qquad
 B_p=\frac{\ii}{c\sqrt{2E_{\mathbf p}}}.
\end{equation*}
对时间求导时，$\partial_t f_p=(\ii E_{\mathbf p}/\hbar)f_p$，所以
\begin{align*}
 \frac{\dd a_t}{\dd t}
 =\int\dd^3 x\,f_p
 \Bigl[&\frac{\ii E_{\mathbf p}}{\hbar}A_p\varphi_H
 +\left(A_p+\frac{\ii E_{\mathbf p}}{\hbar}B_p\right)
 \partial_t\varphi_H
 +B_p\partial_t^2\varphi_H\Bigr].
\end{align*}
中间的一阶时间导数系数为
\begin{align*}
 A_p+\frac{\ii E_{\mathbf p}}{\hbar}B_p
 &=\frac{1}{\hbar c}\sqrt{\frac{E_{\mathbf p}}{2}}
 -\frac{E_{\mathbf p}}{\hbar c\sqrt{2E_{\mathbf p}}}=0.
\end{align*}
余下两项可提出 $B_p$：
\begin{equation*}
 \frac{\dd a_t}{\dd t}
 =\frac{\ii}{c\sqrt{2E_{\mathbf p}}}
 \int\dd^3 x\,f_p
 \left[
 \partial_t^2+\frac{E_{\mathbf p}^2}{\hbar^2}
 \right]\varphi_H.
\end{equation*}
质量壳关系
\begin{equation*}
 E_{\mathbf p}^2=m_{\rm phys}^2c^4+c^2\mathbf p^2
\end{equation*}
把第二项分成质量项与空间动量项。空间部分做两次分部积分；若矩阵元中的波包在空间无穷远衰减，边界项消失，并有
\begin{equation*}
 \int\dd^3 x\,f_p\,(-c^2\nabla^2\varphi_H)
 =\int\dd^3 x\,(-c^2\nabla^2 f_p)\varphi_H
 =\frac{c^2\mathbf p^2}{\hbar^2}
 \int\dd^3 x\,f_p\varphi_H.
\end{equation*}
因此
\begin{equation*}
 \frac{\dd a_t}{\dd t}
 =\frac{\ii c}{\sqrt{2E_{\mathbf p}}}
 \int\dd^3 x\,f_p(x)
 \left(\Box+\frac{m_{\rm phys}^2c^2}{\hbar^2}\right)
 \varphi_H(x).
\end{equation*}
对 $t$ 从 $-\infty$ 积到 $+\infty$，并使用 $\dd^4 x=c\,\dd t\,\dd^3 x$，得到
\begin{equation*}
 \lim_{t\to+\infty}a_t(\mathbf p)
 -\lim_{t\to-\infty}a_t(\mathbf p)
 =\frac{\ii}{\sqrt{2E_{\mathbf p}}}
 \int\dd^4 x\,f_p(x)
 \left(\Box+\frac{m_{\rm phys}^2c^2}{\hbar^2}\right)
 \varphi_H(x).
\end{equation*}
最后用渐近定义~\eqnrefs{eq:asymptotic-annihilator-projection-dp52} 给两端同时乘 $Z^{-1/2}$，即得到\eqnrefs{eq:lsz-one-leg-position-dp52}。

该计算适用于自旋 $0$ 外腿：稳定极点保证 $t\to\pm\infty$ 的单粒子投影存在，空间分部积分要求波包或矩阵元在无穷远没有边界通量，而标量质量壳把 $E_{\mathbf p}^2$ 化成自由标量二次核 $\Box+m_{\rm phys}^2c^2/\hbar^2$。电子外腿对应一阶 Dirac 逆算符，必须使用费米场的 LSZ 约化。
\end{derivation}
